\subsection{THETA3-C13: an explicit algebraic projector onto the Weil space}\label{C13}

\paragraph{Statement.}
Let $A$ be a complex abelian variety of dimension $2n$, with an action of
an imaginary quadratic field $E$ of Weil type: $H^1(A,\Q)$ has rank
$2n$ over $E$, and the two embeddings of $E$ each occur with multiplicity
$n$ on $H^{1,0}(A)$. Put
\[
 W_E(A)=\bigwedge_E^{2n}H^1(A,\Q)\ \subset\ H^{2n}(A,\Q).
\]
Here the inclusion denotes the usual rational two-dimensional Weil
subspace, whose complexification is the sum of the two extreme
eigenspaces. Then the projector onto $W_E(A)$ is represented by an
explicit rational algebraic correspondence on $A\times A$.
Consequently a single algebraic codimension-$n$ cycle $Z$ with
$e_W\cl(Z)\ne0$ suffices to prove that every Weil class of $A$ is
algebraic. More concretely, it suffices to find a smooth projective
$n$-fold $X$, a morphism $b:X\to A$, and $w\in W_E(A)_{\C}$ such that
\begin{equation}\label{C13-period}
 \int_X b^*w\ne0.
\end{equation}
This is an extraction criterion for a supplied cycle, not an assertion
that such a cycle has been constructed for every $A$.

\paragraph{Proof and explicit denominators.}
Let $\mathcal R=E\cap\operatorname{End}(A)$, an order in $E$.
Choose $0\ne b_0\in\mathcal R$ with $\bar b_0=-b_0$, and choose an
integer $M>2n|b_0|$. Set $a=M+b_0$. After choosing one embedding of
$E$ and replacing $b_0$ by its negative if necessary, write
$a=|a|e^{i\vartheta}$ with $0<\vartheta<1/(2n)$.
The $2n+1$ numbers
\[
 \lambda_j=a^j\bar a^{\,2n-j},\qquad 0\le j\le2n,
\]
are distinct: the ratio of two of them is
$e^{2i(j-k)\vartheta}$, whose argument is nonzero and has absolute
value less than $2<2\pi$. Moreover $a^{2n}\notin\R$.
The existence of $b_0$ follows by taking $c-\bar c$ for any nonreal
$c\in\mathcal R$; since $\bar c=\Tr_{E/\Q}(c)-c$ and the trace is an
integer, this difference still belongs to $\mathcal R$.

Let $V_\tau,V_{\bar\tau}$ be the two complex eigenspaces of $H^1(A)$.
The endomorphism $T=[a]^*$ acts on
\[
 H^{2n}(A,\C)=\bigoplus_{j=0}^{2n}
       \bigwedge^jV_\tau\otimes\bigwedge^{2n-j}V_{\bar\tau}
\]
by $\lambda_j$. The two extreme summands are
$W_E(A)_\C$. Define
\[
 f(t)=(t-a^{2n})(t-\bar a^{\,2n}),\qquad
 g(t)=\prod_{j=1}^{2n-1}(t-\lambda_j).
\]
Both are monic in $\Z[t]$: their roots are algebraic integers and
their root sets are stable under the two automorphisms of $E/\Q$.
They are coprime. Put $R=\operatorname{Res}(f,g)\ne0$.
The adjugate of the integral Sylvester matrix supplies
$u,v\in\Z[t]$ with $uf+vg=R$. Thus
\[
 e_W=R^{-1}v(T)g(T)
\]
acts as one on the two extreme summands and as zero on every other
summand of $H^{2n}(A)$. For a projector on the whole cohomology, compose
with the explicit middle K\"unneth projector
\[
 \pi_{2n}=
 \prod_{\substack{0\le k\le4n\\k\ne2n}}
   \frac{[N]^*-N^k\,\mathrm{id}}{N^{2n}-N^k},\qquad N\ge2.
\]
Indeed $[N]^*$ acts by $N^k$ on $H^k(A)$. Endomorphism pullbacks,
their compositions, and the identity are induced by algebraic
correspondences. The displayed polynomial is therefore a rational
algebraic correspondence. Its idempotence is asserted modulo
homological equivalence; no Chow-idempotence is needed here.
An explicit common denominator is $R$ times the product of the
nonzero integers $N^{2n}-N^k$ in the last display.

The two extreme summands have Hodge type $(n,n)$ by the Weil-type
condition. Their rational sum has dimension two, and $T$ on it has
the irreducible quadratic minimal polynomial $f$, because
$a^{2n}\notin\Q$. If $0\ne\alpha=e_W\cl(Z)$, then
$\alpha,T\alpha$ are linearly independent over $\Q$; a dependence
would give a rational eigenvalue of $T$ on this space. Both are
classes of rational algebraic cycles. They therefore span $W_E(A)$.

Finally the middle cup pairing pairs the $j$-summand only with the
$(2n-j)$-summand. In particular $W_E(A)$ is nondegenerate for this
pairing and is orthogonal to the kernel of $e_W$. The projection
formula gives
\[
 \int_X b^*w=\int_A b_*[X]\smile w
           =\int_A e_W\cl(b_*[X])\smile w.
\]
Thus \eqref{C13-period} makes this projected class nonzero, and the
preceding argument applies. $\square$

\framework{This is a concrete projector in the homological category
of abelian motives, followed by algebraic correspondences on its
two-dimensional realization. It gives the explicit form of the
projector-and-rank step C04 in this Weil-type case. The correspondence
framework is Grothendieck, cote~16, scans~43--46, and
\emph{Standard conjectures on algebraic cycles}, pp.~193--196.
It uses no Hodge conjecture to construct the projector. The remaining
condition is a cycle with nonzero projection. In dimension four the
supplied Markman result already establishes algebraicity; this
criterion is not presented as a new proof of that theorem.}

\subsection{THETA3-C14: which curve-product constructions miss the Weil space}\label{C14}

\paragraph{Statement.}
Let $P=C_1\times\cdots\times C_r$ be a product of smooth projective
connected complex curves and let $b:P\to A$ be a morphism to an abelian
variety. The cohomology class $b_*[P]$ belongs to the subalgebra
generated by rational divisor classes on $A$ (and is zero if the
image has dimension less than $r$). If $X$ is smooth projective,
$\pi:X\to P$ is generically finite of degree $d$, and $c=b\circ\pi$, then
$c_*[X]=d\,b_*[P]$ has the same property.

Suppose now $r=n$, $A$ is as in C13, its polarization $\theta$ has
Rosati involution restricting to complex conjugation on $E$, and
$\operatorname{NS}(A)_\Q=\Q\theta$. Then
\[
 e_W\cl(c_*[X])=0.
\]
Thus a construction factoring through a product of $n$ curves cannot
satisfy C13's nonzero-period criterion for this target.

\paragraph{Proof.}
If some $C_i$ has genus zero, every morphism from that factor to an
abelian variety is constant, and the asserted pushforward is zero.
Otherwise choose base points, let $J_i=J(C_i)$, and let
$\iota_i:C_i\to J_i$ be the Abel maps. The universal property of
the Albanese variety expresses $b$ as a translation followed by a
homomorphism
\[
 q:J_1\times\cdots\times J_r\longrightarrow A
\]
composed with $\prod_i\iota_i$. Translation acts trivially on
cohomology. If $g_i=\dim J_i$ and $\theta_i$ is its principal
polarization, the Abel-curve class is
\[
 (\iota_i)_*[C_i]=\frac{\theta_i^{g_i-1}}{(g_i-1)!}.
\]
For completeness, this formula follows by pairing both sides with
$H^2(J_i)$. In a symplectic integral basis $e_j,f_j$ of $H^1(C_i)$,
the Abel map identifies $H^1(J_i)$ with $H^1(C_i)$,
$\theta_i=\sum_j e_j\wedge f_j$, and
$\int_{J_i}\prod_j e_j\wedge f_j=1$. Pairing either displayed class
with $e_j\wedge f_k$ gives $\delta_{jk}$, and pairing with
$e_j\wedge e_k$ or $f_j\wedge f_k$ gives zero. These elements span
$H^2(J_i)$, so Poincar\'e duality proves the formula.

Consequently the class of $P$ in $\prod_iJ_i$ is a product of divisor
classes. Milne, \emph{Lefschetz classes on abelian varieties} (1999),
Corollary~5.5, printed p.~25, states that direct image under a regular
map of abelian varieties preserves such classes. Applying it to $q$
proves the first assertion. The degree formula
$\pi_*[X]=d[P]$ proves the second.

For the final assertion, all codimension-$n$ divisor polynomials on
$A$ are multiples of $\theta^n$. Write
$H^1(A,\C)=V_\tau\oplus V_{\bar\tau}$. The Rosati condition gives
$[a]^*\theta=\Nm_{E/\Q}(a)\theta$ for the element $a$ chosen in C13.
The three eigenvalues on $\bigwedge^2H^1(A)$ are
$a^2,a\bar a,\bar a^2$, which are distinct. It follows that
$\theta\in V_\tau\otimes V_{\bar\tau}\subset\bigwedge^2H^1(A,\C)$.
Thus $\theta^n$ belongs to the $j=n$ summand in C13, whereas the
Weil projector retains only $j=0,2n$. Since $n\ge1$, these summands
are different and $e_W\theta^n=0$. $\square$

\framework{These are Lefschetz classes and algebraic direct images,
in the setting of Milne's Corollaries~5.3 and~5.5, pp.~25--26;
homological and numerical equivalence agree on the Lefschetz
subalgebra there. The Grothendieck anchors are cote~16,
scans~43--46, and the correspondence and positivity framework of
cote~44, scans~7--8. The finding explains the scope of C12's explicit
degree-500 non-CM construction: its diagonal-product cycle is
algebraic, but it supplies no exceptional Weil class. It does not
assert that every map from a finite cover of a curve product factors
through that product. New Albanese factors or a genuinely different
correspondence remain possible.}
